\newpage
\section{习题}

\subsection{基础过关题}

\begin{enumerate}

\item \textbf{（单项选择题）} 下列关于系统总线三组信号线的叙述中，正确的是

\nopagebreak[4]
\begin{center}
\begin{tabular}{ll}
(A) 地址总线是双向的，数据总线是单向的 \\[4pt]
(B) 中断类型号经数据总线传送，总线请求信号经控制总线传送 \\[4pt]
(C) 控制总线中每根信号线的方向相同，均由 CPU 指向 I/O 接口 \\[4pt]
(D) 数据总线的位数决定了主存的最大可寻址空间
\end{tabular}
\end{center}

\ifshowsolutions\par\noindent\textbf{解析：} 逐项判断。(A) 错：地址由主设备发往从设备，只能单向；数据在读写两个方向上流动，是双向的。(B) 对：中断响应时，I/O 接口把中断类型号送到数据线上，CPU 据此查中断向量表；而总线请求 BR、总线允许 BG、中断请求 INTR 都是控制信号，走控制线。(C) 错：控制线每根方向固定，但不同线方向可以不同，INTR、BUSREQ 就是设备发往 CPU 的。(D) 错：最大可寻址空间由地址线位数决定，数据线宽度只决定一次能并行传多少位。\boxed{\textbf{答案：}(B)}\else \answerarea \fi

\item \textbf{（单项选择题）} 采用链式查询方式进行总线仲裁时，下列叙述中\textbf{错误}的是

\nopagebreak[4]
\begin{center}
\begin{tabular}{l}
(A) 只需总线请求 BR、总线忙 BS、总线允许 BG 三根控制信号线 \\[4pt]
(B) 离仲裁器越近的设备优先级越高 \\[4pt]
(C) BG 信号逐级串行传递，某一级链路故障会使其后的设备都得不到总线 \\[4pt]
(D) 可通过改变计数器的初值灵活调整各设备的优先级
\end{tabular}
\end{center}

\ifshowsolutions\par\noindent\textbf{解析：} 链式查询的 BG 信号像接力棒一样从仲裁器出发，按物理位置依次经过各设备，谁先接到谁先得到机会，因此优先级由\textbf{物理位置}固定，最近的最高。它只用 BR、BS、BG 三根线，代价是可靠性差：链路一断，后面的设备全部失联。(D) 描述的是\textbf{计数器定时查询}——那里才有一组设备地址线和可置初值的计数器，通过改变计数起点实现固定、循环或程序可控的优先级。链式查询没有计数器，无法用初值调优先级。\boxed{\textbf{答案：}(D)}\else \answerarea \fi

\item \textbf{（填空题）} 某总线有 16 个主设备。若采用计数器定时查询方式仲裁，除总线请求线 BR 和总线忙线 BS 外，还需要设备地址线，共需控制信号线\underline{\hspace{2cm}}根；若改用独立请求方式，共需控制信号线\underline{\hspace{2cm}}根。

\ifshowsolutions\par\noindent\textbf{解析：} 计数器定时查询要给每个设备一个编号供比较，设备数为 $n=16$ 时需要 $\lceil\log_2 16\rceil=4$ 根地址线，加上 BR、BS 共 $2+4=6$ 根。独立请求方式下每个设备各有一根请求线 BR$_i$ 和一根允许线 BG$_i$，共 $2n=2\times16=32$ 根，且不含公共的 BS。这也是三种集中式仲裁中独立请求连线最多、响应最快的原因。\boxed{\textbf{答案：}\text{6；32}}\else \answerarea \fi

\item \textbf{（填空题）} 某同步总线宽度为 64 位，总线时钟频率为 500 MHz，一个总线周期包含 2 个时钟周期，每个总线周期传送一次数据。该总线的带宽为\underline{\hspace{2cm}}MB/s。

\ifshowsolutions\par\noindent\textbf{解析：} 先把时钟频率换算成\textbf{总线周期}频率，这是本题唯一的坑：一个总线周期占 2 个时钟周期，所以每秒总线周期数 $=500\ \text{MHz}/2=250\times10^{6}$。每次传送 64 位 $=8$ B，故带宽 $=250\times10^{6}\times 8\ \text{B/s}=2000\ \text{MB/s}=2\ \text{GB/s}$。若直接用 $500\times10^{6}\times 8$ 会得到 4000 MB/s，那是把「每时钟周期传一次」和「每总线周期传一次」混为一谈。\boxed{\textbf{答案：}\text{2000（即 2 GB/s）}}\else \answerarea \fi

\item \textbf{（单项选择题）} 异步总线通信中，主设备撤销请求信号必须以收到从设备的应答信号为条件，从设备撤销应答信号也必须以检测到请求信号已撤销为条件。这种握手方式称为

\nopagebreak[4]
\begin{center}
\begin{tabular}{ll}
(A) 不互锁方式 & (B) 半互锁方式 \\[4pt]
(C) 全互锁方式 & (D) 同步方式
\end{tabular}
\end{center}

\ifshowsolutions\par\noindent\textbf{解析：} 三种握手按「谁等谁」区分：不互锁——主设备发完请求，过一段时间就认为传输完成并撤销请求，从设备也不管请求是否还在；半互锁——主设备等从设备的应答后才撤销请求，但从设备发完应答就撤销，不等请求撤销；全互锁——双方互相等待，请求的撤销等应答，应答的撤销等请求撤销，共四次信号变化，最可靠也最慢。题干两句话都是「等对方」，故为全互锁。同步方式根本不用握手信号，靠统一时钟定节拍，不是握手。\boxed{\textbf{答案：}(C)}\else \answerarea \fi

\item \textbf{（填空题）} 某同步总线宽度为 32 位，时钟频率为 100 MHz，每个时钟周期传送一次数据。采用突发传输方式从主存读出 4 个 32 位数据，只需在开始时发送 1 个地址（占 1 个时钟周期），随后连续传送数据，忽略其他开销。完成这次传输需要\underline{\hspace{2cm}}ns，有效数据传输带宽为\underline{\hspace{2cm}}MB/s。

\ifshowsolutions\par\noindent\textbf{解析：} 时钟周期 $T=1/(100\ \text{MHz})=10$ ns。突发传输只需 1 个地址周期，4 个数据各占 1 个周期，共 $1+4=5$ 个时钟周期，时间 $=5\times10=50$ ns。传送的数据量为 $4\times32\ \text{bit}=128\ \text{bit}=16$ B，故有效带宽 $=16\ \text{B}/50\ \text{ns}=320\ \text{MB/s}$。作为对照，该总线的理论峰值为 $4\ \text{B}\times100\ \text{MHz}=400\ \text{MB/s}$，突发传输的效率为 $4/5=80\%$，损失的那 $1/5$ 正是地址周期。\boxed{\textbf{答案：}\text{50；320}}\else \answerarea \fi

\end{enumerate}

\subsection{真题风格与综合训练}

\begin{enumerate}[resume]

\item \textbf{（单项选择题）} 某同步总线采用地址/数据线复用方式，地址线和数据线共 32 根，时钟频率为 200 MHz，每个时钟周期在上升沿和下降沿各传送一次数据。一次突发写事务传送 32 字节数据，其中地址阶段占 1 个时钟周期，数据阶段连续传送，忽略其他协议开销。该事务的有效数据传输带宽约为

\nopagebreak[4]
\begin{center}
\begin{tabular}{ll}
(A) 1.07 GB/s & (B) 1.28 GB/s \\[4pt]
(C) 1.60 GB/s & (D) 2.56 GB/s
\end{tabular}
\end{center}

\ifshowsolutions\par\noindent\textbf{解析：} 分三步算。

（1）数据传送次数：总线宽 32 位 $=4$ B，32 B 需 $32/4=8$ 次传送。

（2）耗时：DDR 使每个时钟周期传 2 次，8 次传送占 $8/2=4$ 个时钟周期；再加 1 个地址周期，共 5 个时钟周期。时钟周期 $T=1/(200\ \text{MHz})=5$ ns，故耗时 $5\times5=25$ ns。

（3）带宽：$32\ \text{B}/25\ \text{ns}=1.28\ \text{GB/s}$。

干扰项来源：(A) 把地址阶段误算成 2 个周期（复用下地址也要一个完整时钟），得 $32/30\approx1.07$ GB/s；(C) 忘记地址周期，用 $32\ \text{B}/20\ \text{ns}=1.6$ GB/s，这恰好等于峰值 $4\ \text{B}\times200\ \text{MHz}\times2=1.6$ GB/s；(D) 把 DDR 的加倍因子重复用了一次。\boxed{\textbf{答案：}(B)}\else \answerarea \fi

\item \textbf{（单项选择题）} 某总线有 8 个主设备采用链式查询仲裁，设备按距仲裁器由近到远编号为 0～7。仲裁器收到请求后产生 BG 信号需 10 ns，BG 每经过一个不请求总线的设备需 5 ns，请求设备接收 BG 不再引入延迟。若只有设备 7 请求总线，从仲裁器收到请求到设备 7 收到 BG 至少需要

\nopagebreak[4]
\begin{center}
\begin{tabular}{ll}
(A) 35 ns & (B) 40 ns \\[4pt]
(C) 45 ns & (D) 50 ns
\end{tabular}
\end{center}

\ifshowsolutions\par\noindent\textbf{解析：} 计数口径：BG 从仲裁器出发，先经过设备 0～6 共 7 个设备，每个转发 5 ns，共 $7\times5=35$ ns，再加上仲裁器自身产生 BG 的 10 ns，共 $45$ ns。（A）只算了 7 级传递、漏掉仲裁器延迟；（B）把级数误算成 6 级（$10+30$）；（D）多算一级，把设备 7 自身的接收也算成一次转发。\boxed{\textbf{答案：}(C)}\else \answerarea \fi

\item \textbf{（综合应用题）} 某同步总线宽度为 128 位，时钟频率为 200 MHz，每个时钟周期传送一次数据。CPU 通过该总线从主存突发读出一个 64 B 的数据块：地址阶段占 1 个时钟周期，随后连续传送 4 个数据，忽略总线申请与仲裁开销。由于主存较慢，从设备在数据阶段开始前通过 WAIT 信号插入了 1 个等待周期。回答下列问题：

\begin{enumerate}
\item[(1)] 该事务共占多少个时钟周期？总时间是多少纳秒？
\item[(2)] 该事务的有效数据传输带宽是多少？总线的理论峰值带宽是多少？带宽利用率为多少？
\item[(3)] 说明 WAIT 等待周期的作用，并解释为什么这种情况下不能靠继续提高总线时钟频率来提升有效带宽。
\end{enumerate}

\ifshowsolutions\par\noindent\textbf{解析：} 时钟周期 $T=1/(200\ \text{MHz})=5$ ns。

(1) 数据个数：$64\ \text{B}/16\ \text{B}=4$ 个。周期数 $=1$（地址）$+1$（等待）$+4$（数据）$=6$ 个，时间 $=6\times5=30$ ns。

(2) 有效带宽 $=64\ \text{B}/30\ \text{ns}=2.13\ \text{GB/s}$（精确值 $32/15$ GB/s）。峰值带宽 $=16\ \text{B}\times200\times10^{6}\ \text{Hz}=3.2\ \text{GB/s}$。利用率 $=2.13/3.2=4/6\approx66.7\%$，即 6 个周期中只有 4 个周期在传有用数据。

(3) WAIT 是半同步通信的手段：总线仍按统一时钟工作，但允许慢速从设备在约定时刻拉低 WAIT，使主设备推迟一个或几个时钟周期再采样数据，从而在不放弃同步总线简单性的前提下兼容慢设备。提高时钟频率会让每个周期变短，从设备物理上更来不及响应，等待周期数反而增加，有效带宽被等待和地址开销摊薄；此时应改用异步握手（由从设备决定何时应答）或分离事务（等待期间把总线让给别的设备），而不是一味提频。\boxed{\textbf{答案：}\text{(1) 6 周期，30 ns；(2) 2.13 GB/s；(3) 见解析}}\else \answerarea[8] \fi

\item \textbf{（综合应用题）} 某系统有 A、B、C 三个主设备，采用集中式独立请求仲裁，固定优先级为 A 最高、B 次之、C 最低，总线不可抢占：一个设备获得总线后必须传完整个事务才能释放。仲裁在事务结束的同一周期内完成，不额外占用总线周期；获得总线的设备从该周期起立即开始传输。各设备的事务长度与请求到达时刻如下（时刻以总线周期为单位，从 0 开始计）：

\nopagebreak[4]
\begin{center}
\begin{tabular}{cccc}
\toprule
设备 & 请求到达时刻 & 事务占用周期数 & 优先级 \\
\midrule
A & 0 & 4 & 高 \\
C & 1 & 2 & 低 \\
B & 2 & 6 & 中 \\
\bottomrule
\end{tabular}
\end{center}

回答下列问题：

\begin{enumerate}
\item[(1)] 按固定优先级调度，写出三个设备获得总线的顺序、各自的开始与完成时刻，以及各自的等待时间（从请求到达至开始传输的周期数）。
\item[(2)] 计算 $[0,12)$ 区间内的总线利用率，并求三个设备的平均等待时间。
\item[(3)] 若把 C 的优先级提到 B 之上（仍低于 A），重新给出服务顺序与平均等待时间，说明优先级设置对平均等待的影响。
\item[(4)] 若 A 以很短的间隔不断发起长事务，会发生什么现象？给出两种在保持独立请求连线的前提下缓解该问题的措施。
\end{enumerate}

\ifshowsolutions\par\noindent\textbf{解析：} 计数约定：设备在周期 $t$ 获得总线，就占用周期 $t,t+1,\dots$，占用 $L$ 个周期后于时刻 $t+L$ 完成并释放；等待时间 $=$ 开始时刻 $-$ 请求时刻。

(1) $t=0$ 只有 A 请求，A 立即获得总线，占用周期 $0\sim3$，于 $t=4$ 完成。$t=4$ 时 B（$t=2$ 到达）和 C（$t=1$ 到达）都在等待，按优先级选 B，B 占用周期 $4\sim9$，于 $t=10$ 完成。之后 C 获得总线，占用周期 $10\sim11$，于 $t=12$ 完成。故顺序为 A$\to$B$\to$C；开始时刻 $0,4,10$；完成时刻 $4,10,12$；等待时间分别为 $0$、$4-2=2$、$10-1=9$ 个周期。

(2) $[0,12)$ 内共 12 个周期，总线被占用 $4+6+2=12$ 个周期，无空闲，利用率为 $12/12=100\%$。平均等待时间 $=(0+2+9)/3=11/3\approx3.67$ 个周期。

(3) C 的优先级高于 B 后：A 仍占 $0\sim3$，$t=4$ 时 C 先获得总线，占 $4\sim5$；B 于 $t=6$ 开始，占 $6\sim11$，$t=12$ 完成。等待时间：A 为 0，B 为 $6-2=4$，C 为 $4-1=3$，平均 $=(0+4+3)/3=7/3\approx2.33$ 个周期。总完成时刻仍是 12，说明优先级不改变总工作量，只改变等待在各设备间的分配；本题中 C 事务短，让它先做使短事务不被长事务挡住，平均等待下降。

(4) A 不断请求且优先级最高时，B、C 每次仲裁都竞争不过 A，可能长期得不到总线，即\textbf{饥饿}。缓解措施（任写两种）：其一，采用循环优先级，让刚服务过的设备降到最低优先级，各设备机会轮转；其二，限制单次占用时间，如规定最长突发长度，超长事务拆成多次并重新仲裁；其三，改用计数器定时查询，通过计数器初值实现循环或程序可控优先级。\boxed{\textbf{答案：}\text{A$\to$B$\to$C；$11/3$；改后 $7/3$；A 独占致饥饿}}\else \answerarea[10] \fi

\end{enumerate}
